\title{Beyond $1+\sqrt2$ for Precedence-Constrained Scheduling of Unit Jobs}
\author{}
\date{}
\maketitle

\begin{abstract}
We study the problem of minimizing total weighted completion time when
unit-time jobs with precedence constraints are scheduled on identical
parallel machines. The number of machines is part of the input. We give
a deterministic polynomial-time approximation algorithm with ratio
$1+\sqrt2-0.0288468<2.385367$. The algorithm adapts Li's
common-quantile rounding framework for a time-indexed linear program.
Our analysis relates the loss in the scalar inequality underlying the
$1+\sqrt2$ bound to progress forced by a band of jobs with nearby
priorities. A geometric estimate charges both failures to enter this
band and movement of the target quantile, and a signed-moment identity
couples the resulting losses. We reduce the verification of the chosen
parameters to finitely many exact quadratic inequalities. The algorithm
schedules only the original jobs and is derandomized by enumerating
polynomially many priority orders.
\end{abstract}

\section{Introduction}
\label{sec:foundations-introduction}

Suppose a collection of tasks must be processed on several identical
machines. Some tasks depend on others: for example, a computation may
require data produced by an earlier computation. Each task takes one unit
of time, a machine can process at most one task at a time, and a task can
start only after all tasks on which it depends have finished. Finishing
some tasks early is more valuable than finishing others early. We express
this preference by giving each task a nonnegative weight and minimizing
the sum of its weight times its completion time. Even with equal
processing times, a decision to run one available task can delay other
tasks and the work that depends on them. The difficulty is to balance
these priorities while respecting both machine capacity and precedence.

Formally, an instance consists of $n$ jobs numbered $1,\ldots,n$,
$m$ identical machines,
nonnegative rational weights $w_j$, and a directed acyclic graph whose
arcs specify precedence constraints. An arc is a directed edge, and acyclic
means that the graph has no directed cycle. Both $n$ and $m$ are positive integers.
Every job starts at or after time zero and occupies one machine for
one unit of time without interruption. If there is an arc from $i$ to
$j$, job $j$ can start only after job $i$ completes. We seek a feasible
schedule minimizing $\sum_{j=1}^{n} w_j\cdot C_j$, where $C_j$ is the completion time
of job $j$. In standard scheduling notation, this is
$P\mid\mathrm{prec},p_j=1\mid\sum_{j=1}^{n}w_j C_j$. The number of machines is
part of the input, and we impose no restriction on the precedence graph.
Here $P$ denotes identical parallel machines, $\mathrm{prec}$ denotes the
precedence constraints, and $p_j$ is the processing time of job $j$.

An \emph{approximation algorithm with ratio} $\rho$ always returns a feasible
schedule of cost at most $\rho$ times the optimum. Our main result is the
following. Its proof appears in Theorem~\ref{thm:final}.

\begin{theorem}[Main result]
\label{thm:main}
There is a deterministic polynomial-time
$(1+\sqrt2-0.0288468)$-approximation for this problem. The ratio is
approximately $2.3853667624$.
\end{theorem}

\section{Related work}

Hall, Schulz, Shmoys, and Wein gave a $(3-1/m)$-approximation for unit
jobs on $m$ machines and a 2-approximation for one machine with arbitrary
processing times \cite[Theorem~2.3 and page~530]{HSSW1997}.
Munier, Queyranne, and Schulz gave a 4-approximation with arbitrary processing
times even when precedence arcs impose nonnegative additional waiting times
\cite{MQS1998}. For arbitrary processing times and every fixed
$\varepsilon>0$, Li improved the ratio to $2+2\ln 2+\varepsilon$
\cite[Theorem~1.1 of the full version]{Li2017}. J\"ager obtained the same
ratio when jobs have release dates and precedence arcs can impose nonnegative
waiting times \cite{Jaeger2018}. A release date is a time before which a job
cannot start.

Earlier scheduling algorithms used a preemptive or fractional schedule to
guide the order of jobs \cite{PhillipsEtAl1998,HSSW1997}. One such priority
is an \emph{$\alpha$-point}, the time by which a fraction $\alpha$ of a job
has been processed in the guiding schedule. Randomized choices of these points
have also been used in approximation algorithms
\cite{ChekuriEtAl2001,SchulzSkutella1997}.

For unit jobs, Li gave a polynomial-time $(1+\sqrt2)$-approximation
with an arbitrary number of machines \cite[Theorem~1.2 and Section~3 of the full version]{Li2017}.
We use the full version's numbering for Li's 2017 paper throughout. The journal
version is \cite{Li2020}.
The analysis uses a time-indexed linear programming relaxation, in which
a job's completion time is represented fractionally, and a common random
quantile to set the priorities of jobs. For every fixed $\varepsilon>0$, a subsequent result of Li gives a
$(1+\sqrt2+\varepsilon)$-approximation in nearly-linear time,
meaning linear in the number of jobs and input precedence arcs up to
factors polylogarithmic in the input size
\cite[Theorem~1.5]{Li2023}. These guarantees apply when the machine count
is part of the input.

Several related results concern different machine or job models. On one
machine with arbitrary processing times, release dates, and precedence
constraints, Sitters and Yang gave a $(2+\varepsilon)$-approximation for
every fixed $\varepsilon>0$ \cite{SittersYang2018}. J\"ager and Warode gave
a 2-approximation for one machine and a 3-approximation for identical
parallel machines in the preemptive setting \cite{JaegerWarode2024}.
Preemption allows a job to be interrupted and resumed later. Their
parallel-machine result applies to arbitrary processing times and runs in
strongly polynomial time.

For \emph{related machines}, each machine has a speed and the time needed to
process a job is its processing requirement divided by that speed.
For arbitrary processing times and precedence constraints, Chudak and Shmoys
gave an $O(\log m)$-approximation for related machines \cite{ChudakShmoys1999}.
Li improved this ratio to $O(\log m/\log\log m)$
\cite[Theorem~1.3 of the full version]{Li2017}. Chekuri, Motwani,
Natarajan, and Stein used single-machine relaxations to design approximation
algorithms for parallel machines and studied special precedence structures,
including trees and series-parallel orders \cite{ChekuriEtAl2001}. On a single
machine, Lawler gave a polynomial-time exact algorithm for series-parallel
orders \cite{Lawler1978}. These orders can be built from individual jobs by
repeatedly placing two disjoint orders in sequence or leaving them unordered
with respect to each other.

Hoogeveen, Schuurman, and Woeginger showed that there is no polynomial-time
approximation scheme for unit jobs with precedence constraints and an
input-dependent number of identical machines, unless $\mathrm{P}=\mathrm{NP}$,
even when all weights are equal \cite{HoogeveenEtAl2001}. Such a scheme would
provide a $(1+\varepsilon)$-approximation for every fixed $\varepsilon>0$.

For every fixed number of machines and fixed $\varepsilon>0$, Garg, Kulkarni, and Li give a
polynomial-time $(2+\varepsilon)$-approximation using a stronger
relaxation \cite[Theorem~1.1]{GKL2019}. Its polynomial-time guarantee is
for fixed $m$ and does not yield a polynomial-time algorithm uniformly
over input-dependent $m$. Under the strengthened Unique Games hypothesis
introduced by Bansal and Khot \cite{BK2009}, approximation within
$2-\varepsilon$ is hard even for one machine \cite[Introduction]{GKL2019}.
Roughly, this hypothesis concerns a regular bipartite graph whose edges
impose permutation constraints on vertex labels. Its vertices are divided
into two sides, every edge joins opposite sides, and every vertex has the
same number of incident edges.
Such a constraint requires one endpoint's label to be a prescribed permutation
of the other endpoint's label.
It asserts hardness of distinguishing instances with a labeling satisfying
every constraint between two nearly complete vertex subsets from instances in which every
labeling satisfies few constraints, while every subset of a specified small
fraction of one side has almost all vertices on the other side as neighbors. The precise quantified
statement is given in \cite[Hypothesis~1]{Svensson2011}. The scheduling hardness
is conditional on this strengthened hypothesis.

\section{Our contribution and proof overview}

Our starting point is Li's linear-programming and common-quantile
framework \cite[Section~3]{Li2017}. The algorithm in this paper uses a
continuous interpolation of the fractional completion distributions and
a machine-driven priority schedule. These adapt Li's framework with a
different choice of rounding rule. For
a chosen quantile, the priority of each job is the time at which its
fractional completion distribution reaches that quantile. The same
quantile is used for every job. The precedence constraints in the linear
program make each predecessor's priority earlier than its successor's.
At each time slot, the algorithm schedules the available jobs of earliest
priority, up to the number of machines.

The basic analysis separates the time until a particular job finishes
into full and underfull slots after restricting the schedule to jobs with
no later priority than the target. To count these jobs at a chosen quantile,
compare each job's fractional completion by the target's fractional time
with that quantile. As the quantile varies, this turns the expected prefix
size into an area problem. A scalar area inequality controls this quantity
and yields the $1+\sqrt2$ baseline. We prove a version of Li's area
inequality and record additional identities for its excess
\cite[Claim~3.3]{Li2017}. Li's analysis already exploits the fact that the workload and chain
bounds cannot simultaneously attain their separate worst cases
\cite[Section~3]{Li2017}. We develop a stronger coupling for this
framework through priority-band progress, geometric excess, and a signed
moment. Li identified an extremal scalar profile for the baseline bound
\cite[Section~3]{Li2017}. In our analysis this profile places its job inside
the priority band for almost every sufficiently large quantile, while the
other equality case has zero pool weight. The proof makes this connection
quantitative.
Specifically, the additional analysis in this manuscript uses
the following ingredients.

\begin{itemize}
\item A priority-band lemma relates idle slots to the number of shallow
jobs (jobs with short prerequisite chains) whose priorities lie in a specified interval. Informally, scheduling
these jobs forces progress on earlier-priority predecessors of the target.
A weighted fractional pool makes this statement useful for the linear
program. We construct the pool so that it automatically has the required size.
\item A geometric estimate controls, together, the jobs that miss the
priority band and the change in the target's quantile over the band. It
charges these losses to the excess in the scalar area inequality.
\item A signed-moment identity couples the charges across jobs instead of
maximizing each contribution independently. This gives a sufficient
parameter certificate for a strict improvement over $1+\sqrt2$. The final
choice is checked using exact inequalities rather than a numerical search
alone.
\end{itemize}

The analysis also uses virtual zero-weight jobs to fill unused fractional
capacity. These jobs are a proof device and are never scheduled by the
algorithm. We prove the transfer needed to apply the padded analysis to
the original instance. Finally, we enumerate the intervals on which all
original jobs have the same priority order. Choosing the least costly
resulting schedule removes randomness in polynomial time.

The fractional relaxation and quantile notation are developed in
Section~\ref{sec:foundations-lp}. Section~\ref{sec:foundations-rounding}
defines the schedule. The scalar inequality, priority-band estimate, and
pool construction follow, after which the geometric and signed-moment
bounds are assembled into a symbolic approximation certificate. The main
text gives the deterministic implementation and resulting guarantee. The
appendix verifies the parameter choice.

